\documentclass[11pt]{article}
\usepackage[margin=1.1in]{geometry}
\usepackage[T1]{fontenc}
\usepackage{lmodern}
\usepackage{amsmath,amssymb,amsthm,mathtools}
\usepackage{microtype}
\usepackage{enumitem}
\usepackage[colorlinks=true,linkcolor=black,citecolor=black,urlcolor=black]{hyperref}

\newtheorem{maintheorem}{Theorem}
\renewcommand{\themaintheorem}{\Alph{maintheorem}}
\newtheorem{lemma}{Lemma}
\theoremstyle{definition}
\newtheorem{definition}{Definition}
\theoremstyle{remark}
\newtheorem*{remark}{Remark}

\newcommand{\R}{\mathbb{R}}
\newcommand{\E}{\mathbb{E}}
\newcommand{\cL}{\mathcal{L}}
\newcommand{\cK}{\mathcal{K}}
\newcommand{\cR}{\mathcal{R}}
\newcommand{\cT}{\mathcal{T}}
\newcommand{\eps}{\varepsilon}
\newcommand{\sm}{\setminus}
\newcommand{\ip}[2]{\langle #1,#2\rangle}
\newcommand{\ND}[1]{\textup{(ND#1)}}

\title{Saddle-to-saddle dynamics of deep polynomial ResNets\\[4pt]\large Statements and proof sketch}
\author{}
\date{November 2026}

\begin{document}
\maketitle

\noindent
This is a compressed version of \texttt{saddle\_to\_saddle.tex}. The statements are the same, and the proofs are reduced to their mechanisms. Constants are not tracked: $O(\cdot)$, $C$ and $c$ are uniform on the bounded sets involved.

\section{Setting}

\paragraph{Network.} Let $\sigma$ be a polynomial with $\sigma(0)=\dots=\sigma^{(k-1)}(0)=0\ne\sigma^{(k)}(0)$ for some $k\ge3$ (e.g.\ $\sigma(z)=z^k$), and
\[
h_0=x,\qquad h_j=h_{j-1}+\sigma(W_jh_{j-1}+b_j)\quad(1\le j\le L),\qquad f_\theta=h_L,
\]
with $W_j\in\R^{d\times d}$ and $b_j\in\R^d$. A \emph{neuron} $n=(j,i)$ has incoming weights $w_n=(W_j[i,:],b_j[i])\in\R^p$, $p=d+1$. These are all the parameters, since the skip connections play the role of outgoing weights: $\theta=(w_n)_{n\in[N]}\in\Theta=(\R^p)^N$, $N=Ld$. Let $F(\theta)=\cL(f_\theta)$, where $\cL$ is a $C^\infty$ loss on the finite-dimensional space of polynomial maps of degree $\le(\deg\sigma)^L$, e.g.\ $\cL(f)=\frac12\E\|f(x)-y\|^2$. Let $\theta^\alpha$ be the gradient flow $\dot\theta=-\nabla F(\theta)$ with $\theta^\alpha(0)=\alpha\xi$, where $\xi\in\Theta$ is fixed and $\alpha\to0$.

\paragraph{Notation.} For a set $A$ of (active) neurons, $S_A=\{\theta:w_n=0\ \forall n\notin A\}$, $\pi_A$ is the orthogonal projection onto $S_A$, $F_A=F|_{S_A}$, $\nabla_A=\pi_A\nabla$, and $\zeta_A(\theta)=\sum_{n\notin A}|w_n|^k$. We write $\theta^{\sm n}$ for $\theta$ with $w_n$ replaced by $0$, and $\iota_n(w)\in\Theta$ for the parameter with $n$-th component $w$ and all others $0$. The \emph{utility} of $n$ at $\theta$ is the $k$-form
\[
U_n(\theta;w)=-\tfrac1{k!}\,D^k_{w_n}F(\theta^{\sm n})[w,\dots,w],\qquad w\in\R^p,
\]
which depends on $\theta$ only through $\theta^{\sm n}$.

\paragraph{Flatness.} Besides smoothness, the proof uses only
\begin{equation}\tag{H}\label{eq:H}
\partial^\beta_{w_n}F=0\ \text{ on }\{w_n=0\}\qquad\text{for every neuron $n$ and every $1\le|\beta|\le k-1$.}
\end{equation}
Indeed, for $n=(j,i)$ we have $f_\theta=\Phi_j\big(\bar h_j+\sigma(\ip{w_n}{\hat h_{j-1}})e_i\big)$, where $\hat h=(h,1)$, $\bar h_j$ is the output of layer $j$ at $w_n=0$, and $\Phi_j$ is the map computed by the layers after $j$. None of these depends on $w_n$. Since $\sigma(t)=t^k\tau(t)$ for a polynomial $\tau$, every monomial of $f_\theta-f_{\theta^{\sm n}}$ in $w_n$ has degree $\ge k$. Hence $F(\theta)-F(\theta^{\sm n})=\sum_{|\beta|=k}w_n^\beta g_\beta(\theta)$ with $g_\beta$ smooth. Depth, width and biases play no further role: everything below holds for every $F\in C^\infty((\R^p)^N)$ satisfying~\eqref{eq:H}.

\begin{lemma}[Flatness]\label{lem:flat}
Locally uniformly in $\theta$, for every neuron $n$ and every set $A\not\ni n$:
\begin{enumerate}[label=\textup{(\alph*)},nosep]
\item $\nabla_{w_n}F(\theta)=-\nabla U_n(\theta;w_n)+O(|w_n|^k)$; in particular $\nabla_{w_n}F(\theta)=O(|w_n|^{k-1})$;
\item $F(\theta)-F(\pi_A\theta)$ and $\nabla_AF(\theta)-\nabla_AF(\pi_A\theta)$ are $O(\zeta_A(\theta))$, $D_{w_n}\nabla_AF(\theta)=O(|w_n|^{k-1})$, and $\nabla U_n(\theta;w)-\nabla U_n(\pi_A\theta;w)=O(|w|^{k-1}\zeta_A(\theta))$.
\end{enumerate}
\end{lemma}

\begin{proof}
Derivatives of $F$ in directions tangent to $\{w_{n'}=0\}$ still satisfy~\eqref{eq:H} in $w_{n'}$. By Taylor's formula in $w_{n'}$, they change by $O(|w_{n'}|^k)$, with a smooth remainder, when $w_{n'}$ is set to $0$. Setting the $w_{n'}$, $n'\notin A$, to zero one at a time gives (b). For (a), by~\eqref{eq:H} the Taylor polynomial of order $k-1$ of $\nabla_{w_n}F$ in $w_n$ at $w_n=0$ is $-\nabla U_n(\theta;w_n)$.
\end{proof}

\paragraph{Consequences.} By (a), $\nabla_{w_n}F=0$ on $\{w_n=0\}$. So every $S_A$ is invariant under the gradient flow, which on $S_A$ is the gradient flow of $F_A$, and critical points of $F_A$ are critical points of $F$. If $w_n=0$, then $F(\theta+\iota_n(w))=F(\theta)-U_n(\theta;w)+O(|w|^{k+1})$. Hence a critical point of $F_A$ at which some $U_n(\theta;\cdot)$, $n\notin A$, takes a positive value is a saddle of $F$. For the mean squared error, the chain rule gives, with $c_\sigma=\sigma^{(k)}(0)/k!$,
\[
U_{(j,i)}(\theta;w)=c_\sigma\,\E\big[\ip{w}{\hat h_{j-1}}^k\,\big(D\Phi_j(h_j)^\top(y-f_\theta)\big)_i\big].
\]
This is the correlation between the $k$-th power of the would-be pre-activation and the back-propagated residual. At $\theta=0$ it equals $c_\sigma\E[\ip{w}{\hat x}^k(y_i-x_i)]$ for every layer $j$.

\section{The limit process and the results}

A \emph{non-degenerate local minimum} of $F_A$ is a point $\theta^*\in S_A$ with $\nabla F_A(\theta^*)=0$ and $\nabla^2F_A(\theta^*)\succ0$; for $A=\emptyset$, $\theta^*=0$. A \emph{non-degenerate local maximum} of a $k$-form $U$ on $S^{p-1}$ is a critical point $u$ of $U|_{S^{p-1}}$ at which the Riemannian Hessian $(\nabla^2U(u)-kU(u)I)|_{u^\perp}$ is negative definite. Let $\theta^*$ be a non-degenerate local minimum of $F_A$, let $n\notin A$, and let $u^*$ be a non-degenerate local maximum of $U_n(\theta^*;\cdot)|_{S^{p-1}}$ with $U_n(\theta^*;u^*)>0$. The \emph{escape orbit} of $\theta^*$ along $(n,u^*)$ is the maximal gradient-flow trajectory $\gamma$ in $S_{A\cup\{n\}}$ with $\gamma(t)\to\theta^*$ and $\gamma_n(t)/|\gamma_n(t)|\to u^*$ as $t\to-\infty$. It exists and is unique up to time translation (Lemma~\ref{lem:escape}).

\begin{definition}[Saddle-to-saddle process]\label{def:process}
Set $s_0=0$, $A_0=\emptyset$, $\theta^{(0)}=0$ and $\Xi_n(0)=\xi_n$. At stage $m$, we are given $s_m$, $A_m$, a non-degenerate local minimum $\theta^{(m)}$ of $F_{A_m}$, and $\Xi_n(s_m)$ for $n\notin A_m$. Each \emph{dormant} $\Xi_n$, $n\notin A_m$, solves
\begin{equation}\label{eq:dormant}
\tfrac{d}{ds}\Xi_n=\nabla U^{(m)}_n(\Xi_n),\qquad U^{(m)}_n:=U_n(\theta^{(m)};\cdot),
\end{equation}
up to its blow-up time $s_m+T^{(m)}_n\in(s_m,\infty]$. Let $s_{m+1}=s_m+\min_{n\notin A_m}T^{(m)}_n$. Stage $m$ is \emph{non-degenerate} if $s_{m+1}<\infty$ and
\begin{enumerate}[label=\textup{(ND\arabic*)},leftmargin=4em,nosep]
\item the minimum is attained at a single neuron $n_{m+1}$;
\item as $s\uparrow s_{m+1}$, $\Xi_{n_{m+1}}(s)/|\Xi_{n_{m+1}}(s)|$ converges to a non-degenerate local maximum $u_{m+1}$ of $U^{(m)}_{n_{m+1}}|_{S^{p-1}}$;
\item the escape orbit $\gamma_{m+1}$ of $\theta^{(m)}$ along $(n_{m+1},u_{m+1})$ is defined for all $t\ge0$ and converges to a non-degenerate local minimum $\theta^{(m+1)}$ of $F_{A_{m+1}}$, where $A_{m+1}=A_m\cup\{n_{m+1}\}$.
\end{enumerate}
In that case, set $\Xi_n(s_{m+1})=\lim_{s\uparrow s_{m+1}}\Xi_n(s)$ for $n\notin A_{m+1}$ and continue. Otherwise stop, and let $M$ be the last stage reached; $M\le N$, since every jump activates a new neuron. The \emph{limit path} is $\bar\theta(s)=\theta^{(m)}$ for $s\in[s_m,s_{m+1})$, $m\le M$. Let $\bar s_n=s_m$ if $n=n_m$, and $\bar s_n=s_M+T^{(M)}_n$ if $n\notin A_M$.
\end{definition}

\begin{remark}
Write $\Xi_n=R_nu_n$ with $|u_n|=1$, and let $U=U^{(m)}_n$. By Euler's identity $\ip{u}{\nabla U(u)}=kU(u)$, \eqref{eq:dormant} reads $\frac{d}{ds}R_n^{2-k}=-k(k-2)\,U(u_n)$ and $\frac{d}{ds}u_n=R_n^{k-2}\,\nabla_{S}U(u_n)$, where $\nabla_S$ is the spherical gradient. So $n$ activates when its accumulated utility $\int_0^sU^{(m(s'))}_n(u_n)\,ds'$ reaches $|\xi_n|^{2-k}/(k(k-2))$. Meanwhile, its direction performs spherical gradient ascent of the utility. This is the alternation of utility maximization and loss minimization of alternating gradient flows (Kunin et al., NeurIPS 2025). In the time $d\tau=R_n^{k-2}ds$, a blow-up forces $\tau\to\infty$ and $U(u_n)\uparrow c>0$. Hence $U^{(m)}_{n_{m+1}}(u_{m+1})>0$ automatically in~\ND2.
\end{remark}

\begin{maintheorem}\label{thm:main}
Let $k\ge3$ and $\xi\in\Theta$. Assume that stages $0,\dots,M-1$ are non-degenerate, and let $S\in(s_M,s_{M+1})$. For $\alpha$ small, $\theta^\alpha$ is defined on $[0,\alpha^{2-k}S]$. As $\alpha\to0$, in the time $s=\alpha^{k-2}t$:
\begin{enumerate}[label=\textup{(\roman*)},nosep]
\item \emph{(plateaus)} for every $\delta>0$, $\theta^\alpha(\alpha^{2-k}s)\to\bar\theta(s)$ uniformly on $[0,S]\sm\bigcup_{1\le m\le M}(s_m-\delta,s_m+\delta)$. In particular the loss converges to the step function $F(\bar\theta(s))$, and $F(\theta^{(0)})>\dots>F(\theta^{(M)})$;
\item \emph{(dormant neurons)} for every neuron $n$ and every $\delta>0$, the rescaled weights $\alpha^{-1}w^\alpha_n(\alpha^{2-k}s)$ converge to $\Xi_n(s)$ uniformly on $[0,\min(S,\bar s_n-\delta)]$;
\item \emph{(jumps)} for $1\le m\le M$, there are times $t^\alpha_m$ with $\alpha^{k-2}t^\alpha_m\to s_m$ such that $\theta^\alpha(t^\alpha_m+\cdot)$ converges to a time translate of $\gamma_m$, locally uniformly on $\R$. Moreover, for $\delta$ smaller than half the gaps between $0,s_1,\dots,s_M,S$, the curve $\{\theta^\alpha(\alpha^{2-k}s):|s-s_m|\le\delta\}$ converges in Hausdorff distance to the closure of $\gamma_m(\R)$.
\end{enumerate}
Each $\theta^{(m)}$ is a critical point of $F$, and a saddle point for $m<M$.
\end{maintheorem}

Let $\cK$ be the smallest set of pairs $(A,\theta)$ with the following properties. First, $(\emptyset,0)\in\cK$. Second, let $(A,\theta)\in\cK$, $n\notin A$, and let $u$ be a non-degenerate local maximum of $U_n(\theta;\cdot)|_{S^{p-1}}$ with $U_n(\theta;u)>0$. If the escape orbit of $\theta$ along $(n,u)$ converges to a non-degenerate local minimum $\theta'$ of $F_{A\cup\{n\}}$, then $(A\cup\{n\},\theta')\in\cK$. The set $\cK$ depends only on the data. Condition~(G) requires, for all $(A,\theta)\in\cK$ and $n\notin A$:
\begin{enumerate}[label=\textup{(G\arabic*)},leftmargin=4em,nosep]
\item every such escape orbit is defined for all $t\ge0$ and converges to a non-degenerate local minimum of $F_{A\cup\{n\}}$;
\item $U_n(\theta;\cdot)|_{S^{p-1}}$ is a Morse function or vanishes identically.
\end{enumerate}
Under (G2), $\cK$ is finite.

\begin{maintheorem}\label{thm:gaussian}
Let $k\ge3$, assume \textup{(G)}, and let $\xi$ have a density, e.g.\ $\xi\sim N(0,I)$. Almost surely, every stage is non-degenerate and Theorem~\ref{thm:main} holds for every $S\notin\{s_1,s_2,\dots\}$. For each $s\ge0$, $\theta^\alpha(\alpha^{2-k}s)\to\bar\theta(s)$ almost surely, and $\bar\theta(s)$ takes values in the finite set $\{\theta:(A,\theta)\in\cK\}$.
\end{maintheorem}

So the data determine the possible plateaus. The initialization only selects which neurons activate, in which order, along which directions, and when.

\section{Proof sketch}

\emph{Idea.} Near a saddle $\theta^*\in S_A$, the active coordinates relax exponentially fast to $\theta^*$. Each inactive weight $w_n=O(\alpha)$ evolves on the slow time scale $\alpha^{2-k}$ by $\dot w_n\approx\nabla U_n(\theta^*;w_n)$ (Lemma~\ref{lem:plateau}). Once a neuron has grown to size $\gg\alpha$ along a non-degenerate maximizer of its utility, it escapes in a vanishing rescaled time while shadowing the escape orbit. The escape is degenerate, since the linearization in $w_n$ vanishes, but in the radial time $\log|w_n|$ it becomes a contraction (Lemma~\ref{lem:escape}). Theorem~\ref{thm:main} follows by induction on the stages, and Theorem~\ref{thm:gaussian} from the scaling of blow-up times.

\emph{Role of $k\ge3$.} It separates the time scales, since $\alpha^{2-k}\gg1$, and it accelerates the relaxation of the active coordinates during the escape, since $\rho^{2-k}\gg1$. It also makes the dormant flows blow up in finite time, at times that depend on $\xi_n$. This breaks the tie between the $L$ neurons of a column, which share the same utility at $\theta=0$. For $k=2$, the dormant flow is linear, a whole column escapes at once, and the outcome of the first jump depends on the ratios of the initial weights.

Throughout, $\theta^*$ is a non-degenerate local minimum of $F_A$, and $\mu>0$ is such that $\nabla^2F_A\succeq\mu$ on $S_A$ near $\theta^*$.

\begin{lemma}[Plateau]\label{lem:plateau}
There are $r_0,C_0>0$ such that the following holds. Let $\Xi_n$, $n\notin A$, solve $\Xi_n'=\nabla U_n(\theta^*;\Xi_n)$ on $[\sigma_0,\sigma_1]$. Let $\tau_\alpha$ be times with $|\pi_A\theta^\alpha(\tau_\alpha)-\theta^*|\le r_0$ and $\alpha^{-1}w^\alpha_n(\tau_\alpha)\to\Xi_n(\sigma_0)$ for $n\notin A$, and let $t(s)=\tau_\alpha+\alpha^{2-k}(s-\sigma_0)$. Then, as $\alpha\to0$, $\alpha^{-1}w_n^\alpha(t(s))\to\Xi_n(s)$ uniformly on $[\sigma_0,\sigma_1]$ for $n\notin A$. Moreover, for small $\alpha$ and $t\in[\tau_\alpha,t(\sigma_1)]$,
\[
|\pi_A\theta^\alpha(t)-\theta^*|\le e^{-\mu(t-\tau_\alpha)}r_0+C_0\sup_{[\tau_\alpha,t]}\zeta_A(\theta^\alpha)=e^{-\mu(t-\tau_\alpha)}r_0+O(\alpha^k).
\]
\end{lemma}

\begin{proof}
While the rescaled weights $\alpha^{-1}w^\alpha_n$, $n\notin A$, stay bounded, $\zeta_A=O(\alpha^k)$. By Lemma~\ref{lem:flat}(b), $e=\pi_A\theta-\theta^*$ satisfies $\dot e=-\nabla F_A(\theta^*+e)+O(\zeta_A)$, and strong convexity and Grönwall's lemma give the bound. In the time $s$, Lemma~\ref{lem:flat} and the homogeneity of $U_n$ give
\[
\tfrac{d}{ds}\big(\alpha^{-1}w_n\big)=-\alpha^{1-k}\nabla_{w_n}F(\theta)=\nabla U_n\big(\theta^*;\alpha^{-1}w_n\big)+O\big(|\pi_A\theta-\theta^*|+\alpha\big).
\]
By the bound, the error term has integral $O(\alpha^{k-2}+\alpha)$ over $[\sigma_0,\sigma_1]$. Grönwall's lemma gives the convergence, and a continuity argument removes the a priori bound.
\end{proof}

\begin{lemma}[Escape]\label{lem:escape}
Let $n\notin A$, $A'=A\cup\{n\}$ and $U=U_n(\theta^*;\cdot)$, and let $u^*$ be a non-degenerate local maximum of $U|_{S^{p-1}}$ with $U(u^*)>0$. Write $\theta=a+\iota_n(\rho u)+z$ with $a=\pi_A\theta$, $\rho=|w_n|$, $|u|=1$, and $z=\theta-\pi_{A'}\theta$, which collects the other inactive neurons.
\begin{enumerate}[label=\textup{(\alph*)},nosep]
\item The escape orbit $\gamma$ of $\theta^*$ along $(n,u^*)$ exists and is unique up to time translation, and $|\pi_A\gamma-\theta^*|=O(|\gamma_n|^k)$. It is the locally uniform limit, as $\eps\to0$, of the gradient-flow trajectories from $\theta^*+\iota_n(\eps u^*)$, shifted so that $|w_n|=\rho_1$ at $t=0$.
\item For every large $K$ there are $\rho_1,r,\eps_1,c,C>0$ with the following property. Suppose that at time $t_0$, $\rho_0:=\rho<\rho_1$, $|u-u^*|\le r$, $\eta:=|z|\le\eps_1\rho_0$ and $|a-\theta^*|\le K(\rho_0^k+\eta^k)$. Then $\rho$ increases and reaches $\rho_1$ at a time $t_1\le t_0+C\rho_0^{2-k}$. On $[t_0,t_1]$, $|z|\le2\eta$, and $\theta(t)$ is $C\big((\rho_0/\rho(t))^c+\eta\big)$-close to the point of $\gamma$ at which $|\gamma_n|=\rho(t)$.
\end{enumerate}
\end{lemma}

\begin{proof}
\emph{Radial time.} By~\eqref{eq:H}, $-\nabla_{w_n}F(a+\iota_n(\rho u)+z)=\rho^{k-1}G(\rho,u,a,z)$, with $G$ smooth for all $u\in\R^p$ and $G(0,u,a,z)=\nabla U_n(a+z;u)$. With $D:=\ip{u}{G}$, the gradient flow reads $\dot\rho=\rho^{k-1}D$, $\dot u=\rho^{k-2}(G-Du)$, $\dot a=-\nabla_AF$. Where $D>0$, we take $\varsigma=\log\rho$ as the time:
\begin{equation}\label{eq:radial}
\frac{du}{d\varsigma}=X_u:=\frac{G-Du}{D},\qquad\frac{da}{d\varsigma}=X_a:=-\rho^{2-k}\,\frac{\nabla_AF}{D},\qquad\frac{dt}{d\varsigma}=\frac{\rho^{2-k}}{D}.
\end{equation}
This system preserves $|u|=1$, since $\frac{d}{d\varsigma}|u|^2=2(1-|u|^2)$. For $z\equiv0$, its solutions give back gradient-flow trajectories in $S_{A'}$.

\emph{Contraction.} At $(\rho,a,z)=(0,\theta^*,0)$, Euler's identity gives $D=kU(u)$ and $X_u=\Phi(u):=\nabla U(u)/(kU(u))-u$. Its zero $u^*$ is a hyperbolic sink. Indeed, since $\nabla^2U(u^*)u^*=(k-1)\nabla U(u^*)$, $D\Phi(u^*)$ is $-2$ on $\R u^*$, and is the Riemannian Hessian of $U|_{S^{p-1}}$ at $u^*$, divided by $kU(u^*)$, on $u^{*\perp}$. The $a$-equation is the gradient flow of the strongly convex $F_A$, accelerated by the factor $\rho^{2-k}/D\to\infty$, plus a forcing term $O(\rho^2+\rho^{2-k}|z|^k)$ (Lemma~\ref{lem:flat}(b)). For $\rho\le\rho_1$ and $\eta\le\eps_1\rho$, consider the regions
\[
\cR_\eta(\rho)=\{(u,a,z):|u-u^*|\le r,\ |a-\theta^*|\le K(\rho^k+\eta^k),\ |z|\le2\eta\}.
\]
On them, $D\ge\kappa>0$, and the Jacobian of $(X_u,X_a)$ with respect to $(u,a)$ has blocks
\[
D_uX_u\preceq-c_0,\qquad D_aX_a\preceq-c\rho^{2-k},\qquad D_aX_u=O(1),\qquad D_uX_a=O(\rho^2).
\]
The last bound holds because $\nabla_AF$ depends on $u$ only through $w_n=\rho u$, and $D_{w_n}\nabla_AF=O(\rho^{k-1})$. Moreover, $(X_u,X_a)$ depends on $z$ only up to $O(|z|)+O\big(\rho^{2-k}(|z|^k+(\rho^k+\eta^k)|z|)\big)=O(\eta)$. Take the norm $\|(\delta u,\delta a)\|^2=|\delta u|^2+\beta^2|\delta a|^2$ with $\beta$ large, and then $\rho_1$ small. The off-diagonal blocks are absorbed, and any two solutions of~\eqref{eq:radial} in these regions, one of them with $z\equiv0$, satisfy
\begin{equation}\label{eq:contraction}
\|y_1-y_2\|(\varsigma)\le e^{-c(\varsigma-\varsigma')}\|y_1-y_2\|(\varsigma')+C\eta,\qquad \varsigma'\le\varsigma\le\log\rho_1.
\end{equation}
The degeneracy helps here: the factor $\rho^{2-k}$ only strengthens the contraction in $a$.

\emph{Invariance.} Under the assumptions of (b), $(u,a,z)(t)\in\cR_\eta(\rho(t))$ on $[t_0,t_1]$. First, $\dot\rho\ge\kappa\rho^{k-1}$ gives $t_1-t_0\le C\rho_0^{2-k}$. Next, $\frac{d}{dt}|z|^{2-k}\ge-C$ by Lemma~\ref{lem:flat}(a), so $|z|^{2-k}\ge\eta^{2-k}-C\rho_0^{2-k}\ge(2\eta)^{2-k}$, because $\eta\le\eps_1\rho_0$ with $\eps_1$ small. On $\{|u-u^*|=r\}$, $X_u$ points inwards, since $\Phi$ does and $X_u-\Phi(u)=O(\rho+|a-\theta^*|+|z|)$ is small. Finally,
\[
\tfrac{d}{d\varsigma}|a-\theta^*|^2\le-\tfrac{2\rho^{2-k}}{D}\,|a-\theta^*|\,\big(\mu|a-\theta^*|-C(\rho^k+\eta^k)\big).
\]
If $\mu K>C$, this is negative on $\{|a-\theta^*|=K(\rho^k+\eta^k)\}$, while this bound increases with $\varsigma$.

\emph{Conclusion.} Let $y^\eps$ solve~\eqref{eq:radial} with $z\equiv0$ and $y^\eps(\log\eps)=(u^*,\theta^*)$. By invariance and~\eqref{eq:contraction} with $\eta=0$, $y^\eps$ converges as $\eps\to0$, uniformly on every $[\varsigma',\log\rho_1]$, to a solution $y^\gamma$ that lies in $\cR_0(e^\varsigma)$ for all $\varsigma$. Moreover $u^\gamma\to u^*$ as $\varsigma\to-\infty$, since $\frac{d}{d\varsigma}|u-u^*|\le-c_0|u-u^*|+O(e^\varsigma)$. As $dt/d\varsigma\ge e^{(2-k)\varsigma}/C$, $t\to-\infty$ when $\varsigma\to-\infty$, and $y^\gamma$ gives $\gamma$. Part~(b) follows from~\eqref{eq:contraction}, applied from $\varsigma'=\log\rho_0$ to the trajectory and to $y^\gamma$, since $|\theta-\gamma|\le|a-a^\gamma|+\rho|u-u^\gamma|+|z|$. For uniqueness, let $\tilde\gamma$ be another trajectory as in~(a). For $t\ll0$, it satisfies $|\pi_A\tilde\gamma-\theta^*|\le K|\tilde\gamma_n|^k$. Otherwise, by the last estimate, this would fail at all earlier times, with $|\pi_A\tilde\gamma-\theta^*|$ increasing backwards in time, which contradicts $\tilde\gamma\to\theta^*$. Then (b) with $\eta=0$ and $t_0\to-\infty$ shows that $\tilde\gamma$ passes through the points of $\gamma$. Hence it is a time translate of $\gamma$.
\end{proof}

\subsection*{Proof of Theorem~\ref{thm:main}}

We prove by induction on $m\le M$:

\smallskip\noindent$(\mathcal C_m)$\quad there are times $\tau^\alpha_m$ with $\alpha^{k-2}\tau^\alpha_m\to s_m$, $|\pi_{A_m}\theta^\alpha(\tau^\alpha_m)-\theta^{(m)}|\le r_0$ (with $r_0$ from Lemma~\ref{lem:plateau} at $\theta^{(m)}$), and $\alpha^{-1}w^\alpha_n(\tau^\alpha_m)\to\Xi_n(s_m)$ for $n\notin A_m$.

\smallskip\noindent$(\mathcal C_0)$ holds with $\tau^\alpha_0=0$.

\emph{Plateau.} Assume $(\mathcal C_m)$. If $m<M$, apply Lemma~\ref{lem:plateau} on $[s_m,s_{m+1}-\delta]$; if $m=M$, apply it on $[s_M,S']$ with $S<S'<s_{M+1}$. This gives (ii) on this interval. Since the inactive weights are $O(\alpha)$, it also gives $|\theta^\alpha-\theta^{(m)}|\le e^{-\mu(t-\tau^\alpha_m)}r_0+O(\alpha)$, which tends to $0$ uniformly for $s\ge s_m+\delta$. This is (i). For $m=0$, (i) holds on the whole interval, since $\pi_\emptyset\theta=0$.

\emph{Jump.} Let $m<M$, $n_*=n_{m+1}$ and $u_*=u_{m+1}$. By \ND2, $|\Xi_{n_*}(s)|\to\infty$ with direction converging to $u_*$ as $s\uparrow s_{m+1}$. Fix $\delta>0$ small, so that $R_\delta:=|\Xi_{n_*}(s_{m+1}-\delta)|$ is large, bounds all $|\Xi_n|$, $n\notin A_m$, on $[s_m,s_{m+1}-\delta]$, and $\Xi_{n_*}(s_{m+1}-\delta)/R_\delta$ is $r/2$-close to $u_*$. Let $t_0$ be the time corresponding to $s_{m+1}-\delta$. For small $\alpha$, the assumptions of Lemma~\ref{lem:escape}(b) at $(A_m,\theta^{(m)},n_*,u_*)$ hold at $t_0$. Indeed, $\rho_0=\alpha(R_\delta+o(1))$, and $\eta=O(\alpha)\le\eps_1\rho_0$ because $R_\delta$ is large. By Lemma~\ref{lem:plateau}, $|a-\theta^{(m)}|\le o(\alpha^k)+C_0N(\alpha R_\delta)^k(1+o(1))\le K\rho_0^k$ if $K\ge2C_0N$. Let $t^\alpha_{m+1}:=t_1$ be the first time after $\tau^\alpha_m$ at which $|w_{n_*}|=\rho_1$; it does not depend on $\delta$. Lemma~\ref{lem:escape}(b) gives:
\begin{itemize}[nosep]
\item $\theta^\alpha(t_1)\to\gamma_{m+1}(0)$, where $\gamma_{m+1}$ is normalized by $|(\gamma_{m+1})_{n_*}(0)|=\rho_1$, since $(\rho_0/\rho_1)^c+\eta\to0$;
\item the jump is fast: $\alpha^{k-2}(t_1-t_0)\le C\alpha^{k-2}\rho_0^{2-k}\to CR_\delta^{2-k}$;
\item $|z|\le2\eta=O(\alpha)$ on $[t_0,t_1]$. So $|\frac{d}{ds}(\alpha^{-1}w_n)|=\alpha^{1-k}|\nabla_{w_n}F|=O(1)$ by Lemma~\ref{lem:flat}(a), and the other dormant neurons move by $O(R_\delta^{2-k})$ in rescaled units.
\end{itemize}
By continuous dependence on initial conditions, forward and backward in time, $\theta^\alpha(t_1+\cdot)\to\gamma_{m+1}$ uniformly on compact intervals. By \ND3, there is $T$ with $|\gamma_{m+1}(T)-\theta^{(m+1)}|\le r_0/2$. Set $\tau^\alpha_{m+1}=t_1+T$. On $[t_1,t_1+T]$, the dormant weights move by $O(\alpha^{k-1})$. Hence $\alpha^{k-2}\tau^\alpha_{m+1}\in[s_{m+1}-\delta,\,s_{m+1}-\delta+CR_\delta^{2-k}]+o(1)$. Moreover, for $n\notin A_{m+1}$,
\[
\big|\alpha^{-1}w_n^\alpha(\tau^\alpha_{m+1})-\Xi_n(s_{m+1})\big|\le CR_\delta^{2-k}+\operatorname{osc}_{[s_{m+1}-\delta,s_{m+1}]}\Xi_n+o(1).
\]
Since $\tau^\alpha_{m+1}$ does not depend on $\delta$, and $R_\delta\to\infty$ as $\delta\to0$, this proves $(\mathcal C_{m+1})$.

\emph{Conclusion.} The plateau and jump steps cover $[0,\alpha^{2-k}S]$, and $\theta^\alpha$ stays bounded there. This gives the existence of $\theta^\alpha$ on this interval, (i), (ii), and the local uniform convergence in (iii). For the Hausdorff convergence at $s_{m+1}$, run the jump step with some $\delta'<\delta$. The curve is close to $\theta^{(m)}$ before $t_0$, to $\gamma_{m+1}([0,T])$ on $[t_1,t_1+T]$, and to $\theta^{(m+1)}$ afterwards. On $[t_0,t_1]$, with $\rho=|w_{n_*}(t)|$, it is $O(\rho+\eta)$-close to $\theta^{(m)}$ and $O((\rho_0/\rho)^c+\eta)$-close to $\gamma_{m+1}$, hence $O(\rho_0^{c/(1+c)}+\eta)$-close to $\overline{\gamma_{m+1}(\R)}$. The converse follows from the local uniform convergence and from (i) at $s_{m+1}\pm\delta$. Finally, $F(\theta^{(m)})>F(\theta^{(m+1)})$ because $F$ decreases strictly along $\gamma_{m+1}$. For $m<M$, $\theta^{(m)}$ is a saddle by the consequences of Lemma~\ref{lem:flat}, since $U^{(m)}_{n_{m+1}}(u_{m+1})>0$.\qed

\subsection*{Proof of Theorem~\ref{thm:gaussian}}

\begin{lemma}[Null sets]\label{lem:null}
Let $U$ be a $k$-form on $\R^p$, $k\ge3$, such that $U|_{S^{p-1}}$ is Morse or vanishes identically. Let $\cT(\zeta)\in(0,\infty]$ be the blow-up time of $\Xi'=\nabla U(\Xi)$, $\Xi(0)=\zeta$. Then $\{\cT=t\}$ is null for every $t<\infty$. Moreover, for almost every $\zeta$ with $\cT(\zeta)<\infty$, $\Xi/|\Xi|$ converges at blow-up to a non-degenerate local maximum of $U|_{S^{p-1}}$. The same holds for piecewise autonomous flows, driven by $\nabla U_j$ on $[\sigma_j,\sigma_{j+1})$.
\end{lemma}

\begin{proof}
By homogeneity, $\cT(c\zeta)=c^{2-k}\cT(\zeta)$ for $c>0$. So the Borel set $\{\cT=t\}$ meets every ray in at most one point, and it is null by polar coordinates. By the Remark, the limit direction is the limit of the spherical gradient flow of $U$ started at $\zeta/|\zeta|$. It fails to be a local maximum only on the stable manifolds of the other critical points, which have positive codimension in $S^{p-1}$. The piecewise case reduces to this one through the flow maps of the earlier pieces, which are diffeomorphisms onto open sets.
\end{proof}

Fix $n$ and $\xi_{-n}=(\xi_{n'})_{n'\ne n}$. A dormant state influences only its own blow-up time. So, until $n$ activates, the process coincides with that of $F|_{\{w_n=0\}}$, which depends on $\xi_{-n}$ only. Moreover, $\Xi_n$ solves a piecewise autonomous flow as in Lemma~\ref{lem:null}, determined by $\xi_{-n}$ and admissible by (G2), started at $\xi_n$. If \ND1 fails at stage $m$ with $n$ among the minimizers, then $\cT(\xi_n)$ equals a number determined by $\xi_{-n}$. If \ND2 fails with $n=n_{m+1}$, then $\xi_n$ lies in the exceptional set of Lemma~\ref{lem:null}. In both cases, the set of such $\xi_n$ is null. By Fubini's theorem, the set of such $\xi$ is null too. Measurability follows from the finiteness of $\cK$ and the lower semicontinuity of blow-up times. \ND3 follows from (G1), since every configuration reached lies in $\cK$. The same argument shows that a fixed $s$ is almost surely not a jump time.\qed

\end{document}
